% Figura vetorial da arquitetura metodologica.
% O arquivo e incluído por artigo_estrutura_inicial.tex dentro de um resizebox.
\begin{tikzpicture}[
  font=\sffamily\scriptsize,
  node distance=4.5mm and 5.5mm,
  >=Latex,
  fluxo/.style={->, line width=0.75pt, draw=black!65},
  fluxoaux/.style={->, line width=0.65pt, draw=black!45, dashed},
  bloco/.style={
    draw=black!55,
    line width=0.65pt,
    rounded corners=2pt,
    align=center,
    inner sep=4.5pt,
    minimum height=10mm
  },
  entrada/.style={bloco, fill=blue!8},
  processo/.style={bloco, fill=black!3},
  aprendido/.style={bloco, fill=orange!13, draw=orange!65!black},
  saida/.style={bloco, fill=green!11, draw=green!55!black},
  familia/.style={
    bloco,
    text width=26mm,
    minimum height=15mm,
    inner sep=3.5pt
  },
  rotulo/.style={font=\sffamily\bfseries\scriptsize, text=black!65},
  nota/.style={
    draw=black!38,
    dashed,
    rounded corners=2pt,
    fill=black!1,
    align=left,
    inner sep=4.5pt
  }
]

% Entrada e decomposicao do video.
\node[entrada, text width=25mm] (video) {\textbf{Vídeo de entrada}\\$V=\{I_1,\ldots,I_T\}$};
\node[processo, text width=37mm, right=of video] (amostragem) {
  \textbf{Amostragem em tempo físico}\\
  timestamps, padronização e controle de qualidade
};
\node[processo, text width=44mm, right=of amostragem] (regioes) {
  \textbf{Regiões semânticas}\\
  rosto, olhos, boca, corpo e fundo\\
  \textcolor{black!62}{track, confiança, origem e fallback}
};

\draw[fluxo] (video) -- (amostragem);
\draw[fluxo] (amostragem) -- (regioes);

% Banco de descritores determinísticos.
\node[familia, fill=blue!9, below=12mm of video, xshift=-2mm] (ga) {
  \textbf{A · Textura e bordas}\\
  LBP, Sobel e Laplaciano
};
\node[familia, fill=cyan!10, right=3.2mm of ga] (gb) {
  \textbf{B · Estrutura local}\\
  SIFT e similaridade de patches
};
\node[familia, fill=violet!9, right=3.2mm of gb] (gc) {
  \textbf{C · Resíduo}\\
  alta passagem bilateral
};
\node[familia, fill=orange!11, right=3.2mm of gc] (gd) {
  \textbf{D · Frequência}\\
  potência e geometria espectral
};
\node[familia, fill=green!10, right=3.2mm of gd] (ge) {
  \textbf{E · Fotometria}\\
  luminância, cor e iluminação
};

\coordinate (familias-top) at ($(gc.north)+(0,4.5mm)$);
\draw[fluxo] (regioes.south) |- (familias-top);
\draw[line width=0.75pt, draw=black!65] (ga.north |- familias-top) -- (ge.north |- familias-top);
\foreach \n in {ga,gb,gc,gd,ge}{
  \draw[fluxo] (\n.north |- familias-top) -- (\n.north);
}

\node[processo, text width=44mm, below=11mm of gb, xshift=-5mm] (sinais) {
  \textbf{Sinais regionais por quadro}\\
  $x_{V,t,r,k}=\phi_k(I_{V,t},r)$\\
  \textcolor{black!62}{inclui contrastes entre regiões}
};
\node[processo, text width=54mm, right=7mm of sinais] (agregacao) {
  \textbf{Agregação estática e temporal}\\
  média, desvio e mediana; $\Delta x$, $\Delta^2x$,\\
  autocorrelação e espectro temporal
};
\node[entrada, text width=24mm, right=7mm of agregacao] (vetor) {
  \textbf{Vetor por vídeo}\\
  $z_V=\Psi(\Phi(I_{1:T}))$
};

\coordinate (familias-bottom) at ($(gc.south)+(0,-4.5mm)$);
\draw[line width=0.75pt, draw=black!65] (ga.south |- familias-bottom) -- (ge.south |- familias-bottom);
\foreach \n in {ga,gb,gc,gd,ge}{
  \draw[fluxo] (\n.south) -- (\n.south |- familias-bottom);
}
\draw[fluxo] (familias-bottom) -| (sinais.north);
\draw[fluxo] (sinais) -- (agregacao);
\draw[fluxo] (agregacao) -- (vetor);

% Etapa aprendida. O fluxo retorna da direita para a esquerda para evitar
% cruzamentos entre as duas faixas da figura.
\node[aprendido, text width=39mm, below=13mm of vetor] (ajuste) {
  \textbf{Ajuste dentro de cada fold}\\
  imputação, escala, seleção e calibração
};
\node[aprendido, text width=43mm, left=7mm of ajuste] (classificador) {
  \textbf{Classificador tabular $g_\theta$}\\
  regressão logística, Random Forest\\
  ou gradient boosting
};
\node[saida, text width=28mm, left=7mm of classificador] (predicao) {
  \textbf{Predição}\\
  $\widehat p(y=\mathrm{fake}\mid V)$
};

\draw[fluxo] (vetor) -- (ajuste);
\draw[fluxo] (ajuste) -- (classificador);
\draw[fluxo] (classificador) -- (predicao);

% Protocolo e auditoria: influencia a avaliacao, mas nao e entrada forense.
\node[nota, text width=137mm, below=11mm of classificador] (protocolo) {
  \textbf{Protocolo de avaliação e auditoria (fora de $z_V$):}
  separação por conteúdo + gerador não visto (LOGOS); baseline de metadados;
  ablações por família, região e temporalidade; controles de movimento, FPS,
  resolução, codec e compressão; holdout comercial após congelamento.
};

\draw[fluxoaux] (protocolo.north -| ajuste.south) -- (ajuste.south);
\draw[fluxoaux] (protocolo.north -| predicao.south) -- (predicao.south);

% Contornos que explicitam o que e fixo e o que e aprendido.
\node[
  draw=blue!55!black,
  line width=0.75pt,
  dashed,
  rounded corners=3pt,
  fit=(amostragem)(regioes)(ga)(ge)(sinais)(agregacao)(vetor),
  inner sep=4.5mm,
  label={[rotulo, text=blue!55!black]above left:extração determinística e auditável}
] (deterministico) {};

\node[
  draw=orange!70!black,
  line width=0.75pt,
  dashed,
  rounded corners=3pt,
  fit=(ajuste)(classificador)(predicao),
  inner sep=4.5mm,
  label={[rotulo, text=orange!70!black]above left:etapa supervisionada}
] (supervisionado) {};

\end{tikzpicture}
